% Rúbrica: Conclusiones (sólidas, basadas en los resultados, con los hallazgos clave y su relevancia).
% Colab no es una limitación: los dos notebooks corrieron ahí el 2026-09-30 sin errores ni avisos
% (docs/evidencia/colab/), y va en el anexo.
\section{Conclusiones}\label{sec:conclusiones}

\subsection{Limitaciones}\label{sec:limitaciones}

\begin{itemize}
  \item \textbf{Es una proxy.} La señal es una reseña negativa temprana, no el arrepentimiento ni el
  reembolso: deja fuera a quien no escribe y a quien escribe después de las dos horas
  (sección~\ref{sec:objetivo}).
  \item \textbf{La muestra retrata el presente.} El \cifraDesdeJulioEntrenamientoPorResena{} de las
  reseñas de entrenamiento es de julio de 2025 en adelante, así que la señal de un juego viejo viene
  de quien lo compra tarde (sección~\ref{sec:datos-periodo}).
  \item \textbf{Hay pocos juegos.} El modelo aprende de \cifraJuegosEntrenamiento{}, y el resultado
  cambia mucho de un fold a otro. Con solo \cifraGratisEntrenamiento{} juegos gratis, el modelo
  extrapola en los gratis, y los \cifraPagoSinPrecioEntrenamiento{} juegos de pago sin precio
  entran con precio 0; los dos casos llevan aviso en la ficha.
  \item \textbf{El orden por score también lo arrastra el precio 0.} En el estante de riesgo bajo de
  Acción, los primeros lugares son de un juego sin precio y de dos gratis, todos con aviso; la
  primera alternativa sin aviso es la tercera tarjeta (sección~\ref{sec:estrategia-recorrido}).
  \item \textbf{Los motivos son palabras clave en inglés.} Solo el \cifraCoberturaMotivosPorResena{}
  de las \cifraNegativasTempranasLimpias{} negativas tempranas del conjunto limpio menciona alguno
  (sección~\ref{sec:interpretabilidad-motivos}).
  \item \textbf{Representatividad.} Las reseñas son en inglés y las escriben compradores de todo el
  mundo, mientras que los precios son de la tienda de México y el usuario objetivo compra desde
  México. Tampoco hay datos de PlayStation ni de Xbox.
  \item \textbf{No hubo un estudio con usuarios externos.} La usabilidad se apoya en decisiones de
  diseño y en pruebas automáticas, no en sesiones con personas que no conocían el sitio.
  \item \textbf{La infraestructura es gratuita.} La API se duerme cuando nadie la usa y su disco
  borra las calificaciones y los comentarios al reiniciarse (sección~\ref{sec:arquitectura}).
\end{itemize}

\subsection{Qué se concluye}\label{sec:conclusiones-hallazgos}

La señal de arrepentimiento temprano depende sobre todo del juego: la varianza de su tasa entre
juegos es unas \cifraSobredispersion{} veces la esperada si todos tuvieran la misma tasa, y lo que
mejor la acompaña es la crítica. Un modelo pequeño, una regresión logística con cinco variables que se conocen antes
de comprar, ordena los juegos que no vio mejor que el trivial: \cifraPRAUCCociente{} veces en la
validación agrupada y \cifraPRAUCExternoCociente{} veces en la prueba externa, con un intervalo que
no baja de \cifraExternoICInf{}. Es una ventaja real y modesta, y el documento la reporta siempre
junto al trivial.

El perfil de quien compra no la mejora más allá del ruido entre folds, y el dato que lo
representaba en el sitio no estaba validado. Por eso el riesgo es del título y vale igual para
cualquier persona. El aporte del proyecto es convertir esa ventaja en algo que sirve antes de
pagar: una banda con palabras y no una probabilidad, factores que dicen qué tan firme es su
evidencia, motivos con su cobertura a la vista y acciones concretas, como comparar, buscar una
alternativa o aprovechar la ventana de reembolso. Lo que no hace es decirle a nadie qué tan probable
es que se arrepienta, ni explicar con certeza por qué un juego decepciona.

\subsection{Siguientes pasos}\label{sec:siguientes}

La Parte A, los modelos que leen el texto de las reseñas y la comparación entre algoritmos, sigue en
la rama \texttt{modelos-texto} con su prerregistro, y entrará a este documento cuando llegue a
master (sección~\ref{sec:modelacion-texto}). La Parte B parte de la cobertura de los motivos: probar
si un modelo que lee el texto completo asigna motivos a más negativas tempranas que las palabras
clave, contra un conjunto de referencia etiquetado a mano y con la regla de decisión fijada antes de
mirar. Más allá de eso, el modelo ganaría con más juegos, en especial gratis, y con una prueba de la
interfaz con personas que no la conocen.
